\subsection{Barriers of the IS in Porto Marghera}
Although the industrial symbiosis of Porto Marghera seems at first to be a sustainable and robust system, several barriers still limit the development of the ecodistrict.\\

\noindent The main barrier is the central role of Veritas, which treats and distributes waste from the municipality of Venice to other companies. Almost all material flows within the industrial symbiosis pass through Veritas. This creates a risk for the resilience and sustainability of the system. Even if Veritas is owned by the City of Venice and does not seek economic profit, any major incident for the company that stops its activities would affect the entire industrial symbiosis. This would directly impact the production and activities of all companies involved. As a result, this lack of resilience could discourage new companies from joining the Porto Marghera ecodistrict.\\

\noindent Other barriers that may affect the sustainability of this case study have also been identified. The close geographical location of the companies, which reinforces the resilience issue, could become a problem in the event of a disaster or a local disruption (fire, flooding, strike, etc.). Another barrier concerns the treatment of mixed waste. The quality of the mixed waste must be high enough to be directly used by companies. If the quality is too low, the waste must be reprocessed by Veritas, which leads to additional costs and delays.\\
\subsection{Possible solutions}
To reduce the impact of these barriers, several solutions can be considered to make the industrial symbiosis of Porto Marghera more sustainable, resilient, and reliable. One solution is to involve new companies in order to diversify the exchanged material flows. For example, the activities of Veritas could be supported by another waste treatment company, or by companies that could replace key functions in case of failure. This would increase the overall resilience of the system. Another solution is to duplicate Veritas’ treatment and distribution lines to ensure reliable input supplies, even if one line is affected by an incident.
Despite the significant reuse of recycled waste, the industrial symbiosis of Porto Marghera still faces several barriers that reduce its resilience and the reliability of its supply chains. Therefore, improving the sustainability of the system is necessary to attract new companies and to further develop industrial symbiosis activities.